% !TEX program = pdflatex
% Build: pdflatex -interaction=nonstopmode condensed_coefficient_sections_EN.tex (twice)
\documentclass[11pt,a4paper]{article}
\usepackage[a4paper,margin=24.5mm,headheight=15pt]{geometry}
\usepackage[T1]{fontenc}
\usepackage[utf8]{inputenc}
\usepackage{lmodern}
\usepackage{amsmath,amssymb,amsthm,mathtools,mathrsfs}
\usepackage{tikz-cd}
\usepackage{booktabs,array,enumitem}
\usepackage{microtype}
\usepackage{xcolor}
\usepackage{fancyhdr}
\usepackage[unicode,colorlinks=true,linkcolor=blue!50!black,citecolor=blue!50!black,urlcolor=blue!50!black]{hyperref}
\usepackage{bookmark}
\setlength{\parskip}{2.5pt}
\setlength{\parindent}{1.25em}
\setlist[enumerate]{leftmargin=2.3em,itemsep=2pt,topsep=3pt}
\setlist[itemize]{leftmargin=2.3em,itemsep=2pt,topsep=3pt}
\numberwithin{equation}{section}
\newtheorem{theorem}{Theorem}[section]
\newtheorem{proposition}[theorem]{Proposition}
\newtheorem{lemma}[theorem]{Lemma}
\newtheorem{corollary}[theorem]{Corollary}
\theoremstyle{definition}
\newtheorem{definition}[theorem]{Definition}
\newtheorem{example}[theorem]{Example}
\theoremstyle{remark}
\newtheorem{remark}[theorem]{Remark}
\newcommand{\ProFin}{\mathrm{ProFin}}
\newcommand{\Cond}{\mathrm{Cond}}
\newcommand{\Sh}{\mathrm{Sh}}
\newcommand{\CAlg}{\mathrm{CAlg}}
\newcommand{\Mod}{\mathrm{Mod}}
\newcommand{\CRing}{\mathrm{CRing}}
\newcommand{\Sch}{\mathrm{Sch}}
\newcommand{\Ring}{\mathrm{Ring}}
\newcommand{\D}{\mathrm D}
\newcommand{\LC}{\mathrm{LC}}
\newcommand{\Hom}{\mathrm{Hom}}
\newcommand{\Map}{\mathrm{Map}}
\newcommand{\id}{\mathrm{id}}
\newcommand{\proet}{\mathrm{pro\acute{e}t}}
\newcommand{\OO}{\mathcal O}
\newcommand{\FF}{\mathcal F}
\newcommand{\MM}{\mathcal M}
\newcommand{\I}{\mathcal I}
\newcommand{\ringunder}[1]{\underline{#1}}
\newcommand{\Qlbar}{\overline{\mathbb Q}_{\ell}}
\newcommand{\Zlbar}{\overline{\mathbb Z}_{\ell}}
\newcommand{\xo}[1]{\left(#1\right)}
\pagestyle{fancy}
\fancyhf{}
\fancyhead[L]{\footnotesize Sectionwise condensed coefficient rings}
\fancyhead[R]{\footnotesize \thepage}
\renewcommand{\headrulewidth}{0.3pt}
\title{\bfseries Sectionwise Condensed Coefficient Rings\\[3pt]
  and Their Pullback to the Pro-\'{e}tale Site\\[5pt]
  \large A pullback-oriented lax Cartesian coefficient system}
\author{Technical construction notes}
\date{October 2026}
\begin{document}
\maketitle
\vspace{-1.5em}
\begin{abstract}
Fix a prime $\ell$.  All coefficient rings in these notes are defined
\emph{by their sections on profinite test spaces}, rather than by limits,
colimits, or algebraic localizations of sheaves.  For the four local
coefficient rings $\mathbb Z_\ell$, $\mathbb Q_\ell$, $\mathcal O_E$ and
$E$, the defining sections are continuous functions $C(S,T)$ with the
usual local-field topologies.  For $\overline{\mathbb Q}_\ell$, a section
is instead an equivalence class of a finite extension $E/\mathbb Q_\ell$
and a continuous map $S\to E$.  We verify the sheaf axioms directly and
only \emph{afterwards} record the familiar limit, localization, and
filtered-colimit descriptions as comparison isomorphisms.  We construct
the canonical geometric morphism $p_X:X_{\proet}\to *_{\proet}$ from
$S\mapsto X[S]$, describe the pulled-back ring sheaves explicitly,
and exhibit their coherent structure maps
$f^{-1}A_Y\xrightarrow{\sim}A_X$.  The orientation is fixed as part
of the definition of a pullback-oriented lax Cartesian coefficient
system.  No six-operation formalism is developed.
\end{abstract}
\begingroup\small\setlength{\parskip}{0pt}\tableofcontents\endgroup

\section{Precise scope and the meaning of ``lax Cartesian''}

We work in a fixed size universe, or equivalently with a suitably large
regular cardinal so that the pro-\'{e}tale sites used below are essentially
small.  Let $\mathscr S$ be a category of schemes closed under the base
changes required below.  We use ordinary sheaves of sets, abelian groups,
and commutative rings; there is no need for solid modules or analytic rings.

For $X\in\mathscr S$, write $X_{\proet}$ for the pro-\'{e}tale site of
Bhatt--Scholze: its objects are weakly \'{e}tale morphisms $U\to X$ and
its coverings are fpqc covering families.  The notation
$\Sh(X_{\proet})$ refers to the associated topos.

\begin{definition}[The direction used throughout]\label{def:direction}
Let $\mathscr R(X)$ denote the category of commutative ring objects in
$\Sh(X_{\proet})$.  For $f:X\to Y$, let
$f^{-1}:\mathscr R(Y)\to\mathscr R(X)$ be the inverse-image functor.
A \emph{pullback-oriented lax Cartesian ring system} is the following data:
\begin{enumerate}
\item an object $A_X\in\mathscr R(X)$ for every $X$;
\item for every $f:X\to Y$, a ring morphism
\begin{equation}\label{eq:lax-direction}
     \boxed{\quad\theta_f:f^{-1}A_Y\longrightarrow A_X\quad};
\end{equation}
\item unit and composition coherences, made explicit in
  Definition~\ref{def:lax-formal}.
\end{enumerate}
If every $\theta_f$ is invertible, we call the system
\emph{Cartesian} (in this pullback-oriented convention).
\end{definition}

\begin{remark}[Convention warning: lax versus oplax]\label{rem:convention}
The direction in \eqref{eq:lax-direction} is intentional.  For a
contravariant pseudofunctor $\mathscr R:\mathscr S^{\mathrm{op}}\to\mathsf{Cat}$,
the usual \emph{Grothendieck Cartesian fibration} over $\mathscr S$ has
arrows over $f:X\to Y$ corresponding to maps
$A_X\to f^{-1}A_Y$ in $\mathscr R(X)$---the reverse direction.
Our convention can be expressed using fiberwise opposite categories, or as
one chosen orientation of a lax/oplax natural transformation.
Authors do not uniformly attach the word ``lax'' to the same orientation.
Here we \emph{define} the convention by \eqref{eq:lax-direction}.
The distinction disappears for the coefficient systems constructed below,
since all their structural arrows are isomorphisms.  Also, this use of
``lax Cartesian'' is \emph{not} the separate notion of a lax Cartesian
\emph{symmetric monoidal structure} on a category-valued functor.
\end{remark}

\section{A condensed ring, constructed by its sections}

\subsection{The pro-\'{e}tale site of a point}

Let $\ProFin$ be the category of profinite sets and continuous maps.
Equip it with the Grothendieck topology in which
$\{S_i\to S\}_{i=1}^r$ is a covering family precisely when
$\coprod_i S_i\to S$ is surjective.  Define
\[
    *_{\proet}:=(\ProFin,J),\qquad
    \Cond(\mathrm{Set})=\Sh(*_{\proet}),\qquad
    \Cond(\mathrm{Ab})=\Sh(*_{\proet},\mathrm{Ab}).
\]
A \emph{condensed commutative ring} is simply a sheaf of commutative rings
on this site.  Equivalently, it is a commutative unital ring object of
$\Cond(\mathrm{Set})$, or a commutative algebra object of
$\Cond(\mathrm{Ab})$ with its sheaf tensor product.

The following explicit construction will be used for every coefficient.

\begin{definition}[Continuous-function condensed ring]\label{def:cond-T}
Let $T$ be a Hausdorff topological commutative ring.  Define a presheaf
of rings $\ringunder T$ on $\ProFin$ by
\begin{equation}\label{eq:point-sections}
      \boxed{\quad \ringunder T(S):=C(S,T)
       =\{a:S\to T\mid a\text{ is continuous}\}.\quad}
\end{equation}
For a continuous map $u:S'\to S$, the restriction map is
\[
      u^*:C(S,T)\longrightarrow C(S',T),\qquad a\longmapsto a\circ u.
\]
The ring operations, constants, and unit are defined pointwise:
\[
(a+b)(s)=a(s)+b(s),\quad
(ab)(s)=a(s)b(s),\quad 1_S(s)=1_T.
\]
\end{definition}

\begin{proposition}[The sheaf axiom, without abstraction]\label{prop:sheaf}
The presheaf $S\mapsto C(S,T)$ is a sheaf of commutative rings on
$*_{\proet}$.  A continuous ring homomorphism $T\to T'$ gives a morphism
of condensed rings $\ringunder T\to\ringunder{T'}$ by postcomposition.
\end{proposition}
\begin{proof}
A finite jointly surjective family is equivalent, for the sheaf axiom, to
a surjection $q:S'\twoheadrightarrow S$ of profinite sets.  Since $S'$ is
compact and $S$ is Hausdorff, $q$ is a topological quotient map.  Thus a
continuous $a:S'\to T$ descends to a continuous $S\to T$ if and only if
$a$ is constant on the fibers of $q$, i.e., if and only if
$\operatorname{pr}_1^*a=\operatorname{pr}_2^*a$ on $S'\times_S S'$.
Consequently
\[
 C(S,T)\xrightarrow{\sim}
 \operatorname{Eq}\bigl(
 C(S',T)\rightrightarrows C(S'\times_S S',T)\bigr).
\]
This is an equalizer of rings because all operations are pointwise.
Postcomposition clearly respects restrictions, descent, and ring operations.
\end{proof}

\begin{example}[Sections with finite values]\label{ex:finite-test}
If $A$ is a finite discrete commutative ring and $S$ is profinite, a
section of $\ringunder A$ is a locally constant function $a:S\to A$.
Equivalently, there is a finite quotient $q:S\twoheadrightarrow S_0$
and a function $a_0:S_0\to A$ with $a=a_0\circ q$.  Addition and
multiplication simply combine these finite tables of values on a common
finite quotient.  This description does not invoke a categorical
limit or colimit.
\end{example}

\begin{remark}[The methodological rule]
We distinguish \emph{definitions of coefficients} from \emph{identities
subsequently satisfied by those coefficients}.  Every definition in the
next section specifies a ring of actual sections, restriction by
precomposition, and pointwise or finite-stage ring operations.  Some
limits, colimits, and localization identities will be proved later;
none of them is used to define a coefficient ring.
\end{remark}

\section{Five coefficient rings defined by their sections}

Fix a prime $\ell$.  We work inside a fixed algebraic closure
$\Qlbar$ of $\mathbb Q_\ell$, so that all finite extensions below are
actual subfields of one field.  We use the usual $\ell$-adic topology
on $\mathbb Z_\ell$ and $\mathbb Q_\ell$ and the usual local-field
topologies on $\OO_E$ and $E$.

\subsection{\texorpdfstring{The condensed $\mathbb Z_\ell$}{The condensed Zl}: continuous integral sections}

\begin{definition}\label{def:Zl}
For each profinite set $S$, \emph{define}
\begin{equation}\label{eq:Zl-def}
  \boxed{\quad \ringunder{\mathbb Z_\ell}(S):=C(S,\mathbb Z_\ell).\quad}
\end{equation}
For a continuous map $u:S'\to S$, define $u^*a=a\circ u$.
Addition, multiplication, zero, and one are defined pointwise.  Thus
$\ringunder{\mathbb Z_\ell}$ is a condensed commutative ring by
Proposition~\ref{prop:sheaf}.
\end{definition}

Here is a fully explicit alternative description of a \emph{section},
not a second definition of the ring.  To $a\in C(S,\mathbb Z_\ell)$
associate its residue functions
\[
   a_n:S\longrightarrow\mathbb Z/\ell^n\mathbb Z,
   \qquad a_n(s)=a(s)\bmod\ell^n.
\]
Every $a_n$ is locally constant and
$a_{n+1}(s)\bmod\ell^n=a_n(s)$.  Conversely, any family of locally
constant functions $(a_n)_{n\geq1}$ satisfying these pointwise
congruences specifies a unique continuous $\mathbb Z_\ell$-valued
function: at each $s$ the residues specify an $\ell$-adic integer, and
continuity follows because the residue maps define the topology.
Operations are computed residue by residue:
\[
 (a+b)_n=a_n+b_n,\qquad (ab)_n=a_nb_n.
\]
Notice that the \emph{object defined} in \eqref{eq:Zl-def} is the
continuous-function ring, not the inverse limit of a tower of sheaves.

\subsection{\texorpdfstring{The condensed $\mathbb Q_\ell$}{The condensed Ql}: continuous rational sections}

\begin{definition}\label{def:Ql}
For every profinite $S$, \emph{define}
\begin{equation}\label{eq:Ql-def}
   \boxed{\quad\ringunder{\mathbb Q_\ell}(S):=C(S,\mathbb Q_\ell).\quad}
\end{equation}
Again restrictions are given by precomposition and ring operations
are pointwise.  The field carries its ordinary nonarchimedean topology.
\end{definition}

\begin{lemma}[Uniform denominators]\label{lem:uniform-denominator}
A function $a:S\to\mathbb Q_\ell$ belongs to
$\ringunder{\mathbb Q_\ell}(S)$ if and only if there exist an integer
$N\geq0$ and $b\in C(S,\mathbb Z_\ell)$ such that
\[
                   a(s)=\ell^{-N}b(s)\quad(s\in S).
\]
In particular, the integer $N$ is \emph{uniform in $s$}.
\end{lemma}
\begin{proof}
The image of a continuous function on a profinite set is compact in
$\mathbb Q_\ell$, hence bounded.  Therefore some $\ell^N a(S)$ is
contained in $\mathbb Z_\ell$, and multiplication by $\ell^N$ is
continuous.  The converse is immediate.
\end{proof}

Thus rational sections admit a second \emph{concrete representation}
by pairs $(N,b)$ with $b:S\to\mathbb Z_\ell$ continuous, where
\[
  (N,b)\sim(M,c)\quad\Longleftrightarrow\quad
                     \ell^M b=\ell^N c
\]
as functions on $S$.  Sums and products are calculated after choosing
a common denominator; for example, with $K=\max(N,M)$,
\[
 (N,b)+(M,c)=\bigl(K,\ell^{K-N}b+\ell^{K-M}c\bigr),\qquad
 (N,b)(M,c)=(N+M,bc).
\]
The pairs are merely a presentation of the ring already defined by
continuous functions.

\subsection{\texorpdfstring{The condensed $\OO_E$ and $E$}{The condensed OE and E}: integral and rational sections}

Let $E/\mathbb Q_\ell$ be a finite extension, let $\OO_E$ be its ring
of integers, and choose a uniformizer $\varpi_E$.

\begin{definition}\label{def:OE-E}
For each profinite set $S$, \emph{define} two rings by
\begin{equation}\label{eq:OE-E-def}
  \boxed{\quad
   \ringunder{\OO_E}(S):=C(S,\OO_E),\qquad
   \ringunder E(S):=C(S,E).
  \quad}
\end{equation}
Restriction is precomposition and arithmetic is pointwise in the
target topological ring.  Both are condensed commutative rings.
\end{definition}

An integral section $a\in C(S,\OO_E)$ is concretely determined by its
locally constant residues
\[
 a_n:S\to\OO_E/\varpi_E^n,\qquad
 a_{n+1}\bmod\varpi_E^n=a_n.
\]
Conversely every compatible family of such residue functions gives a
unique continuous $\OO_E$-valued section.  A rational section
$c\in C(S,E)$ has \emph{one uniform denominator}:
\[
 \exists N\geq0,\ b\in C(S,\OO_E),\qquad c=\varpi_E^{-N}b.
\]
The proof is the same compactness argument as for
Lemma~\ref{lem:uniform-denominator}.  The choice of $\varpi_E$ is
used only in the alternative presentations, not in the definitions
\eqref{eq:OE-E-def}.

If $E\subset F$ are two finite subextensions of $\Qlbar$, the
continuous inclusions yield actual morphisms of condensed rings
\begin{equation}\label{eq:finite-transitions}
 \ringunder{\OO_E}\longrightarrow\ringunder{\OO_F},\qquad
 \ringunder E\longrightarrow\ringunder F,
\end{equation}
whose component on $S$ is simply postcomposition with the inclusion.
No extra comparison isomorphism is needed to compose these maps.

\subsection{\texorpdfstring{The condensed $\overline{\mathbb Q}_\ell$}{The condensed Qbarl}: explicit finite-field data}

The algebraic closure requires a \emph{different primary definition}.
We do \emph{not} begin by assigning to $\Qlbar$ a topology and then
considering all continuous maps $S\to\Qlbar$.

\begin{definition}[Algebraic-closure sections]\label{def:Qbar}
For each profinite $S$, let $\mathscr P(S)$ be the set of pairs
\[
 (E,a),\qquad
 \mathbb Q_\ell\subset E\subset\Qlbar,\quad [E:\mathbb Q_\ell]<\infty,
 \qquad a\in C(S,E).
\]
Declare $(E,a)\sim(F,b)$ if there exists a finite extension
$L\subset\Qlbar$ containing both $E$ and $F$ such that the functions
$S\to L$ induced by $a$ and $b$ are equal.  Define
\begin{equation}\label{eq:Qbar-def}
 \boxed{\quad
  \ringunder{\Qlbar}(S):=\mathscr P(S)/{\sim}.
 \quad}
\end{equation}
For $u:S'\to S$, define $u^*[(E,a)]=[(E,a\circ u)]$.
To define the ring operations, select any finite common extension
$L\supset E,F$ and set
\begin{align}\label{eq:Qbar-arithmetic}
 [(E,a)]+[(F,b)]&:=[(L,a+b)],\\
 [(E,a)]\cdot[(F,b)]&:=[(L,ab)].
\end{align}
The zero and one are the classes of the constant maps to
$0,1\in\mathbb Q_\ell$.  These definitions are independent of the
chosen common finite extension.
\end{definition}

The equivalence relation is especially concrete.  Since all fields
are subfields of the fixed $\Qlbar$, two pairs are equivalent exactly
when
\[
                      a(s)=b(s)\text{ in }\Qlbar
                    \quad\text{for every }s\in S.
\]
As a nontrivial example, the identity map
$\mathbb Z_\ell	o\mathbb Q_\ell$ represents a section on the
profinite space $S=\mathbb Z_\ell$; it is not locally constant and its
image is infinite.
The useful distinction is that \emph{each representative is a continuous
map into one finite local field}; we have not defined a notion of
continuity into the naked set $\Qlbar$ itself.

\begin{proposition}[Direct sheaf verification]\label{prop:Qbar-sheaf}
The assignment $S\mapsto\ringunder{\Qlbar}(S)$, with restrictions
and operations from Definition~\ref{def:Qbar}, is a condensed
commutative ring.  Its underlying ring of sections at a singleton is
$\ringunder{\Qlbar}(*)\cong\Qlbar$.
\end{proposition}
\begin{proof}
A common finite extension exists for any \emph{finite} set of
representatives, by taking their compositum.  Hence
\eqref{eq:Qbar-arithmetic} is well-defined, and the ring axioms can
all be checked inside one finite field containing the finitely many
representatives occurring in any given identity.

For the sheaf condition, consider a finite jointly surjective family
$\{S_i\to S\}_{i=1}^r$ and matching sections
$[(E_i,a_i)]\in\ringunder{\Qlbar}(S_i)$.  Choose a common finite
field $L$ containing all $E_i$.  The $a_i$, regarded as continuous
$L$-valued functions, agree on every overlap
$S_i\times_S S_j$.  Since the surjection
$\coprod_iS_i\twoheadrightarrow S$ is a quotient map, they descend
to a unique continuous function $a:S\to L$.  The class $[(L,a)]$
is the unique glued section.  Thus finite-cover descent holds, and
the empty-object condition is immediate.  Evaluation on a singleton
identifies the classes of pairs $(E,a)$ with the elements of the
algebraic closure.
\end{proof}

\begin{remark}[Why this is not simply $C(S,\Qlbar)$]
The symbol $C(S,\Qlbar)$ does not have a meaning until one chooses a
topology on $\Qlbar$.  The metric topology inherited from
$\mathbb C_\ell$ admits continuous maps whose values do not lie in
any one finite extension.  For example, after choosing primitive
$\ell^n$-th roots of unity $\zeta_{\ell^n}$, the sequence
$\ell^n\zeta_{\ell^n}$ tends to $0$ in the metric topology, although
no finite extension contains all its terms.  Such a function on a
profinite convergent sequence is \emph{not} one of our sections.
One can separately compare our ring to continuous functions for a
carefully specified final topology, but this is not part of its
definition, and we do not need that comparison in what follows.
\end{remark}

\subsection{Limit, localization and colimit identities --- theorems, not definitions}

Now that the five coefficient rings have been defined as actual
section functors, their familiar algebraic descriptions can be
\emph{proved} and may safely be used later.

\begin{proposition}[Comparison identities]\label{prop:comparison-only}
There are natural isomorphisms of condensed commutative rings
\begin{align}
 \ringunder{\mathbb Z_\ell}
   &\simeq\varprojlim_n\ringunder{\mathbb Z/\ell^n\mathbb Z},
      \label{eq:Zl-comparison}\\
 \ringunder{\mathbb Q_\ell}
   &\simeq\ringunder{\mathbb Z_\ell}[1/\ell],
      \label{eq:Ql-comparison}\\
 \ringunder{\OO_E}
   &\simeq\varprojlim_n\ringunder{\OO_E/\varpi_E^n},
      \label{eq:OE-comparison}\\
 \ringunder E
   &\simeq\ringunder{\OO_E}[1/\varpi_E],
      \label{eq:E-comparison}\\
 \ringunder{\Qlbar}
   &\simeq\varinjlim_{E/\mathbb Q_\ell\ \mathrm{finite}}\ringunder E.
      \label{eq:Qbar-comparison}
\end{align}
\emph{None} of these isomorphisms is taken as a definition.
\end{proposition}
\begin{proof}
The first and third identities follow by taking compatible residue
functions: the bijections are natural in the profinite test space and
preserve ring operations.  The second and fourth follow from the
uniform-denominator statements, again naturally in the test space.
For the last identity, the right-hand side has on each profinite test
space exactly the classes of pairs $(E,a)$ specified in
Definition~\ref{def:Qbar}.  This uses that filtered colimits of
sheaves on the profinite point are computed sectionwise: its coverings
are finite, and filtered colimits commute with finite-limit sheaf
conditions.  The comparison map is the identity on the explicit
representatives and respects all ring operations.
\end{proof}

\begin{remark}[What these comparisons do not assert]
The inverse image of a geometric morphism is not claimed to preserve
arbitrary inverse limits.  In the next section, integral inverse-limit
identities \emph{after} pullback will be checked from the concrete
continuous-function description.  Filtered-colimit descriptions after
pullback must be understood as statements about \emph{sheaves}, not
necessarily as sectionwise identities on non-quasi-compact objects.
\end{remark}

\section{\texorpdfstring{The geometric morphism $p_X$ from the profinite point}{The geometric morphism pX from the profinite point}}

\subsection{\texorpdfstring{Explicitly attach a pro-\'{e}tale $X$-scheme to $S$}{Explicitly attach a pro-etale X-scheme to S}}

For a finite set $I$, define the finite \'{e}tale $X$-scheme
\begin{equation}\label{eq:finite-XI}
          X[I]:=\coprod_{i\in I} X.
\end{equation}
If $S=\varprojlim_i S_i$ is a profinite set expressed as an inverse
limit of finite discrete sets, define
\begin{equation}\label{eq:XS}
   \boxed{\quad X[S]:=\varprojlim_i X[S_i]
             =\varprojlim_i\coprod_{s\in S_i}X.\quad}
\end{equation}
The transition morphisms are finite \'{e}tale over $X$ and hence affine;
the inverse limit exists as an $X$-scheme and is pro-\'{e}tale over $X$.
This construction is independent of the chosen presentation of $S$ up
to a unique canonical isomorphism.  A continuous map $u:S\to T$ induces
an $X$-morphism $X[u]:X[S]\to X[T]$.

\begin{lemma}[Properties of profinite test objects]\label{lem:XS}
The functor
\begin{equation}\label{eq:site-functor}
             u_X:\ProFin\longrightarrow X_{\proet},
                      \qquad S\longmapsto X[S]
\end{equation}
preserves finite limits and sends jointly surjective finite covering
families to pro-\'{e}tale coverings.  In particular,
\[
 X[*]\simeq X,\qquad
 X[S\times T]\simeq X[S]\times_X X[T],\qquad
 X[S\sqcup T]\simeq X[S]\sqcup X[T].
\]
\end{lemma}
\begin{proof}
Each statement is immediate for finite sets, where $X[I]$ is a finite
disjoint union of copies of $X$.  The general statements follow by
passing to inverse limits through finite quotients.  Surjectivity on
finite fibers is stable under base change to $X$, so the resulting
families cover in the pro-\'{e}tale topology.
\end{proof}

\subsection{The direction of the site morphism}

\begin{proposition}[The canonical map to the point]\label{prop:pX}
The continuous left-exact functor $u_X$ of \eqref{eq:site-functor}
induces a geometric morphism of topoi
\begin{equation}\label{eq:pX}
     \boxed{\quad
       p_X:\Sh(X_{\proet})\longrightarrow
          \Sh(*_{\proet})=\Cond(\mathrm{Set}).\quad}
\end{equation}
Its inverse-image functor
\[
  p_X^{-1}:\Cond(\mathrm{Set})\longrightarrow\Sh(X_{\proet})
\]
preserves finite limits and all colimits.  On sheaves of abelian groups
it is exact and strong symmetric monoidal; in particular, it takes
condensed commutative rings to commutative ring sheaves on $X_{\proet}$.
\end{proposition}
\begin{proof}
The direction of the geometric morphism is opposite to the direction
of the site functor $u_X$.  Concretely, $p_X^{-1}F$ is obtained by
left Kan extension of the presheaf $F$ along $u_X$ (in the appropriate
opposite categories), followed by sheafification.  The general theory
of morphisms of sites supplies its right adjoint and left exactness.
A geometric inverse image preserves all colimits because it is a left
adjoint; preserving finite limits makes it exact on abelian sheaves.
Since tensor products of abelian sheaves are formed by sheafifying the
presheaf tensor product, the inverse image is strong symmetric monoidal.
\end{proof}

\begin{remark}[Why one must not just precompose]
It would be incorrect to define $p_X^{-1}F$ by writing
$F(U)$ for an arbitrary weakly \'{e}tale $U\to X$: such a $U$ is not an
object of the profinite-point site.  The left Kan extension and
sheafification are essential in the abstract definition; the next
subsection explains why the answer nevertheless has a very simple
continuous-function description for our particular rings.
\end{remark}

\subsection{Compute the inverse image of continuous-function coefficients}

Recall that $X_{\proet}$ has a basis of w-contractible affine objects
$U\to X$.  For such a $U$, the connected-component space $\pi_0(U)$
is profinite; in fact the property of being w-contractible makes
covers of these basis objects split.

\begin{lemma}[Point-pullback formula on the basis]\label{lem:point-eval}
For a condensed set, abelian group, or ring $F$ and a w-contractible
affine $U\in X_{\proet}$, there is a natural identification
\begin{equation}\label{eq:point-eval}
             (p_X^{-1}F)(U)\simeq F(\pi_0(U)).
\end{equation}
\end{lemma}
\begin{proof}[Explanation]
For any profinite set $S$ there is a natural bijection
\[
      \Hom_X(U,X[S])\simeq C(\pi_0(U),S).
\]
For finite $S$ this means that a morphism to a finite disjoint union
of copies of $X$ amounts to a finite clopen decomposition of $U$;
the general case follows by its finite quotients.  In the
inverse-image site computation, these identifications reduce the
value on a w-contractible affine to evaluation of $F$ at the
profinite set $\pi_0(U)$.  The splitting of coverings on this basis
ensures that sheafification does not change that value.  This is the
usual point-to-pro-\'{e}tale inverse-image calculation (compare
\cite[Example~3.1]{HRS} and \cite[Section~4.2]{BS}).
\end{proof}

\begin{proposition}[Continuous-function pullback]\label{prop:continuous-X}
For $T\in\{\mathbb Z_\ell,\mathbb Q_\ell,\OO_E,E\}$, there is a
natural isomorphism of ring sheaves
\begin{equation}\label{eq:continuous-X}
 \boxed{\quad
    p_X^{-1}\ringunder T
        \xrightarrow{\sim}
          \bigl[U\longmapsto C(|U|,T)\bigr]^\#.
 \quad}
\end{equation}
In fact the indicated continuous-function presheaf is already a
sheaf on $X_{\proet}$.
\end{proposition}
\begin{proof}
Each of the four topological rings is Hausdorff and totally
disconnected.  On a w-contractible affine $U$, every continuous
$|U|\to T$ factors uniquely through the profinite connected-component
space $\pi_0(U)$, so
\[
 C(|U|,T)\simeq C(\pi_0(U),T)
   =\ringunder T(\pi_0(U))
   \simeq(p_X^{-1}\ringunder T)(U).
\]
The identifications respect restrictions and ring operations, so they
extend from the basis to all of $X_{\proet}$.  The sheaf property of
$U\mapsto C(|U|,T)$ is \cite[Lemma~4.2.12]{BS}; the inverse-image
identification is also given in \cite[Example~3.1]{HRS}.
\end{proof}

\subsection{What the five pulled-back ring sheaves actually are}

For \emph{any} condensed commutative ring $\Lambda$ already defined on
the point, define its pullback to $X$ by
\begin{equation}\label{eq:LambdaX}
                  \boxed{\quad\Lambda_X:=p_X^{-1}\Lambda.\quad}
\end{equation}
This is a definition by inverse image of a \emph{previously defined}
ring object, not a limit or colimit definition of coefficients.

For the four topological local rings, Proposition~\ref{prop:continuous-X}
gives the concrete formula, on every pro-\'{e}tale object $U$:
\begin{align*}
 \mathbb Z_{\ell,X}(U)&=C(|U|,\mathbb Z_\ell),
 &\mathbb Q_{\ell,X}(U)&=C(|U|,\mathbb Q_\ell),\\
 \OO_{E,X}(U)&=C(|U|,\OO_E),
 &E_X(U)&=C(|U|,E).
\end{align*}
All equalities here refer to the canonical sheaf identifications, not
to an arbitrary choice of presentation.

\begin{proposition}[The algebraic-closure pullback is locally finite-stage]
\label{prop:Qbar-X}
Let
\[
                    \mathcal B_X:=p_X^{-1}\ringunder{\Qlbar}.
\]
For an arbitrary $U\in X_{\proet}$, a section of $\mathcal B_X(U)$
may be described by \emph{local finite-field data}:
\begin{enumerate}
 \item a pro-\'{e}tale covering family $\{U_i\to U\}$;
 \item finite subextensions $E_i/\mathbb Q_\ell$ in $\Qlbar$;
 \item continuous functions $a_i:|U_i|\to E_i$;
 \item agreement of $a_i$ and $a_j$, as functions with values in
       $\Qlbar$, on $U_i\times_U U_j$.
\end{enumerate}
Two such sets of data determine the same section precisely when they
agree after a common refinement of covers.  Addition and
multiplication are performed locally after enlarging the finitely
many fields appearing in a given arithmetic operation on an overlap.
On a w-contractible affine $U$ there is the particularly simple formula
\begin{equation}\label{eq:Qbar-X-basis}
 \boxed{\quad
 \mathcal B_X(U)\simeq
  \{(E,a)\mid E/\mathbb Q_\ell\ \mathrm{finite},\,
                a\in C(|U|,E)\}/{\sim}.
 \quad}
\end{equation}
\end{proposition}
\begin{proof}
The comparison \eqref{eq:Qbar-comparison}, which was \emph{proved after}
constructing the condensed ring, together with preservation of
colimits by $p_X^{-1}$, identifies $\mathcal B_X$ with the sheaf
associated to the presheaf of finite-stage pairs
$U\mapsto\{(E,a:|U|\to E)\}/{\sim}$.  By the definition of
sheafification this means exactly locally defined pairs satisfying
descent on overlaps, modulo refinement.  On a w-contractible affine,
Lemma~\ref{lem:point-eval} gives
$\mathcal B_X(U)\simeq\ringunder{\Qlbar}(\pi_0(U))$.
Continuous functions on $|U|$ with values in any $E$ factor through
$\pi_0(U)$, yielding \eqref{eq:Qbar-X-basis}.
\end{proof}

\begin{remark}[A global section need not have one finite field]
For non-quasi-compact $U$, the fields $E_i$ appearing in the local
representatives of a section can grow without bound over infinitely
many connected components.  Thus it would be incorrect to replace
the locally finite-stage description by one global representative
$(E,a)$ for arbitrary $U$.  Nor do we write
$\mathcal B_X(U)=C(|U|,\Qlbar)$ without an independently specified
topology and comparison theorem.
\end{remark}

\begin{remark}[Identities after pullback, not new definitions]
As consequences of Proposition~\ref{prop:comparison-only} and the
continuous-function calculation, one has canonical \emph{sheaf}
isomorphisms
\[
\begin{gathered}
 \mathbb Z_{\ell,X}\simeq
    \varprojlim_n(\mathbb Z/\ell^n)_X,\qquad
 \OO_{E,X}\simeq
    \varprojlim_n(\OO_E/\varpi_E^n)_X,\\
 \mathbb Q_{\ell,X}\simeq\mathbb Z_{\ell,X}[1/\ell],\qquad
 E_X\simeq\OO_{E,X}[1/\varpi_E],\qquad
 \mathcal B_X\simeq\varinjlim_E E_X.
\end{gathered}
\]
The inverse-limit statements are special properties checked from
continuous functions, not instances of a general inverse-limit
preservation principle for $p_X^{-1}$.  The last filtered colimit is
taken in sheaves and does not require one uniform finite field on
all of a non-quasi-compact $U$.
\end{remark}

\section{Functorial pullback along a map of schemes}

\subsection{A natural identity before taking any sheaves}

Let $f:X\to Y$ be a scheme morphism.  For each finite set $I$,
\[
        X\times_Y\Bigl(\coprod_{i\in I}Y\Bigr)
                   \simeq\coprod_{i\in I}X.
\]
Taking inverse limits over a finite-quotient presentation of a
profinite set $S$ yields a canonical isomorphism
\begin{equation}\label{eq:basechange-XS}
               \boxed{\quad
                X\times_Y Y[S]\xrightarrow{\sim}X[S].\quad}
\end{equation}
This is natural in $S$ and $f$, and compatible with iterated base change.

The scheme morphism $f$ induces a geometric morphism, also written $f$,
\[
   f:\Sh(X_{\proet})\longrightarrow\Sh(Y_{\proet}),
\]
using base change of weakly \'{e}tale objects $V\to Y$ to
$X\times_Y V\to X$.  Formula~\eqref{eq:basechange-XS} provides a
canonical comparison of geometric morphisms
\begin{equation}\label{eq:triangle}
        p_X\ \cong\ p_Y\circ f.
\end{equation}
The comparison is coherent for identity maps and compositions of
scheme morphisms, because it comes directly from ordinary fiber-product
associativity.

\begin{theorem}[The fundamental coefficient pullback]
\label{thm:alpha}
For every condensed commutative ring $\Lambda$ and every $f:X\to Y$
there is a canonical isomorphism of \emph{ring sheaves}
\begin{equation}\label{eq:alpha}
   \boxed{\quad
        \alpha_f^\Lambda:
        f^{-1}\Lambda_Y=f^{-1}p_Y^{-1}\Lambda
                 \xrightarrow{\sim}p_X^{-1}\Lambda=\Lambda_X.\quad}
\end{equation}
These isomorphisms are natural in $\Lambda$, and satisfy the
composition coherence of Theorem~\ref{thm:cart-section}.
\end{theorem}
\begin{proof}
Apply inverse image to the canonical isomorphism of geometric
morphisms \eqref{eq:triangle}.  Both inverse-image functors are exact
strong symmetric monoidal functors on abelian sheaves, so the resulting
comparison is an isomorphism of commutative ring objects.
The coefficient naturality and all coherence conditions follow from
the corresponding statements for \eqref{eq:basechange-XS}.
\end{proof}

\subsection{The comparison on explicit sections}

For $T\in\{\mathbb Z_\ell,\mathbb Q_\ell,\OO_E,E\}$, let
$V\to Y$ be pro-\'{e}tale and set $V_X=X\times_YV$.  A continuous
section $h\in C(|V|,T)$ has the obvious base-changed section
\begin{equation}\label{eq:precomposition-function}
 h\longmapsto h\circ|\operatorname{pr}_V|
                         \in C(|V_X|,T).
\end{equation}
For $\Qlbar$, a local representative $(E,a)$ is similarly sent to
$(E,a\circ|\operatorname{pr}_V|)$; this commutes with the equivalence
relation, refinement, restrictions, and finite-stage arithmetic.
These explicit maps are compatible with the canonical isomorphism
$\alpha_f^\Lambda$ in Theorem~\ref{thm:alpha}.  Thus the direction
$f^{-1}\Lambda_Y\to\Lambda_X$ is both geometrically canonical and
concretely visible on sections.

\section{The pullback-oriented lax Cartesian system}

\subsection{The formal unit and composition axioms}

The categories of ring sheaves, together with geometric inverse images,
form a contravariant pseudofunctor
\begin{equation}\label{eq:R-pseudo}
 \mathscr R:\mathscr S^{\mathrm{op}}\longrightarrow\mathsf{Cat},
 \qquad \mathscr R(X):=\CRing(\Sh(X_{\proet})),
 \qquad \mathscr R(f):=f^{-1}.
\end{equation}
For a composable pair
$X\xrightarrow{f}Y\xrightarrow{g}Z$, write
\[
    c_{f,g}:f^{-1}g^{-1}\xRightarrow{\sim}(g\circ f)^{-1}
\]
for the canonical composition constraint of the pseudofunctor.

\begin{definition}[Coherent lax Cartesian coefficient data]
\label{def:lax-formal}
A pullback-oriented lax Cartesian ring system
$(A_X,\theta_f)$ satisfies:
\begin{enumerate}
\item \emph{Unit:} $\theta_{\id_X}=\id_{A_X}$, after the canonical
  identification $\id_X^{-1}A_X\simeq A_X$;
\item \emph{Composition:} for $X\xrightarrow fY\xrightarrow gZ$,
  the following equality of ring-sheaf morphisms holds:
\begin{equation}\label{eq:coherence}
   \boxed{\quad
   \theta_{g\circ f}\circ c_{f,g,A_Z}
      =\theta_f\circ f^{-1}(\theta_g):
            f^{-1}g^{-1}A_Z\longrightarrow A_X.\quad}
\end{equation}
\end{enumerate}
The same equation is required coherently for triple compositions;
in a strictification it is an equality, and for an ordinary
pseudofunctor it is read together with its standard associativity
constraint.
\end{definition}

\begin{theorem}[The coefficient rings form a Cartesian system]
\label{thm:cart-section}
For any condensed commutative ring $\Lambda$, put
\[
       A_X=p_X^{-1}\Lambda,
       \qquad \theta_f:=\alpha_f^\Lambda.
\]
Then $(A_X,\theta_f)$ is a pullback-oriented lax Cartesian ring system.
In fact every $\theta_f$ is an isomorphism, so it is Cartesian.
Explicitly, for $X\xrightarrow fY\xrightarrow gZ$ the diagram
\begin{equation}\label{eq:coherence-diag}
\begin{tikzcd}[column sep=large,row sep=large]
 f^{-1}g^{-1}\Lambda_Z
    \arrow[r,"f^{-1}(\alpha_g^\Lambda)"]
    \arrow[d,"c_{f,g}"']
  & f^{-1}\Lambda_Y \arrow[d,"\alpha_f^\Lambda"]\\
 (g\circ f)^{-1}\Lambda_Z
    \arrow[r,"\alpha_{g\circ f}^\Lambda"']
  & \Lambda_X
\end{tikzcd}
\end{equation}
commutes, and the identity constraint is the identity.
\end{theorem}
\begin{proof}
The objects and structural maps have been constructed in
\eqref{eq:LambdaX} and \eqref{eq:alpha}.  For an identity map $X\to X$,
the base-change isomorphism $X\times_XX[S]\simeq X[S]$ is the identity.
For a composition $X\to Y\to Z$, both paths in
\eqref{eq:coherence-diag} are induced by the canonical identification
\[
 X\times_Y\bigl(Y\times_Z Z[S]\bigr)
                   \xrightarrow{\sim}X\times_ZZ[S].
\]
The associativity law of fiber products supplies the triple-composition
coherence.  All $\alpha_f^\Lambda$ are isomorphisms by
Theorem~\ref{thm:alpha}.
\end{proof}

\begin{corollary}[Functoriality in the coefficient ring]
\label{cor:coeff-functor}
A morphism $a:\Lambda\to\Lambda'$ of condensed commutative rings
induces a morphism of Cartesian coefficient systems,
\[
       a_X:=p_X^{-1}a:\Lambda_X\longrightarrow\Lambda'_X,
\]
in the sense that for every $f:X\to Y$ the square
\begin{equation}\label{eq:coeff-square}
\begin{tikzcd}[column sep=large,row sep=large]
 f^{-1}\Lambda_Y \arrow[r,"\alpha_f^\Lambda"]
                \arrow[d,"f^{-1}a_Y"']
       & \Lambda_X \arrow[d,"a_X"]\\
 f^{-1}\Lambda'_Y \arrow[r,"\alpha_f^{\Lambda'}"']
       & \Lambda'_X
\end{tikzcd}
\end{equation}
commutes.  In particular, all inclusions and localizations among
$\mathbb Z_\ell$, $\mathbb Q_\ell$, $\OO_E$, $E$, and $\Qlbar$
are compatible with scheme pullback.
\end{corollary}
\begin{proof}
Apply the naturality in $\Lambda$ of the geometric comparison
$p_X\cong p_Y\circ f$.
\end{proof}

\subsection{The associated module systems, explicitly}

It is often useful to formulate the lax structure not only for rings
but also for modules, since it makes the next categorical layer immediate.
For a general pullback-oriented lax ring system $(A_X,\theta_f)$,
let
\[
   \MM_A(X)=\Mod_{A_X}\bigl(\Sh(X_{\proet},\mathrm{Ab})\bigr).
\]
Define the module pullback along $f:X\to Y$ by the familiar formula
\begin{equation}\label{eq:module-pullback}
    f_A^*M
    := A_X\otimes_{f^{-1}A_Y}f^{-1}M,
\end{equation}
where the ring homomorphism giving the algebra structure on $A_X$ is
\emph{exactly} $\theta_f:f^{-1}A_Y\to A_X$.
The coherence~\eqref{eq:coherence} identifies the iterated pullback
$f_A^*g_A^*$ canonically with $(g\circ f)_A^*$.
For general noninvertible $\theta_f$, the tensor product in
\eqref{eq:module-pullback} need not be exact; nothing here asserts
otherwise.

For the coefficients $A_X=\Lambda_X$ pulled back from the point,
$\theta_f$ is an isomorphism.  Therefore the formula simplifies to
\begin{equation}\label{eq:module-simple}
      \boxed{\quad
        f_\Lambda^*M\simeq f^{-1}M,
     \quad}
\end{equation}
with the $\Lambda_X$-action transported along
$\alpha_f^\Lambda$.  It is consequently exact on sheaves of modules,
strong symmetric monoidal, and compatible with compositions.

\begin{definition}[Pullback-oriented lax module section]
A lax module section over a fixed coefficient system $A$ consists of
$M_X\in\MM_A(X)$ for every $X$ and $A_X$-linear morphisms
\[
          \rho_f:f_A^*M_Y\longrightarrow M_X
\]
for every $f:X\to Y$, satisfying the same unit and composition
identities as in Definition~\ref{def:lax-formal}, with the module
pullback comparison in place of $c_{f,g}$.  It is Cartesian when every
$\rho_f$ is invertible.
\end{definition}

\begin{proposition}[Concrete Cartesian sections from the point]
\label{prop:point-module}
Let $N$ be a module over the condensed ring $\Lambda$ on
$*_{\proet}$.  For every $X$, put $N_X=p_X^{-1}N$, with its natural
$\Lambda_X$-module structure.  Then the canonical isomorphisms
\begin{equation}\label{eq:point-module}
       f_\Lambda^*N_Y\xrightarrow{\sim}N_X
\end{equation}
make $(N_X)$ into a Cartesian module section.
In particular, taking $N=\Lambda$ gives the unit section
$X\mapsto\Lambda_X$.
\end{proposition}
\begin{proof}
The isomorphisms~\eqref{eq:point-module} follow from
$f^{-1}p_Y^{-1}N\simeq p_X^{-1}N$ and
\eqref{eq:module-simple}.  Their units and compositions are the same
canonical identifications used in Theorem~\ref{thm:cart-section}.
\end{proof}

\begin{remark}[The category-valued system is a different datum]
The assignment $X\mapsto\MM_\Lambda(X)$ with pullbacks
$f_\Lambda^*$ is a contravariant \emph{pseudofunctor of categories}.
An individual family $(M_X,\rho_f)$ is a lax or Cartesian
\emph{section} of that pseudofunctor.  These should not be conflated.
If one subsequently works with derived or stable categories, the same
coherence of $\alpha_f^\Lambda$ supplies the underlying geometric
pullback pseudofunctor; no exceptional functor is needed for this
coefficient-level construction.
\end{remark}

\section{Minimal implementation checklist}

The constructions can be reused in the following order.  The first
stage \emph{only} uses concrete sections; comparison isomorphisms are
subsequent results.
\begin{enumerate}
 \item \textbf{Four topological coefficient functors.}
 $\ringunder T(S)=C(S,T)$ for
 $T=\mathbb Z_\ell,\mathbb Q_\ell,\OO_E,E$;
 restriction is precomposition and operations are pointwise.
 \item \textbf{Define the algebraic-closure coefficient functor:}
 $\ringunder{\Qlbar}(S)$ consists of equivalence classes $[(E,a)]$
 where $a:S\to E$ is continuous and $E$ is finite over
 $\mathbb Q_\ell$; implement arithmetic in a common finite extension.
 \item \textbf{Prove the sheaf conditions directly:}
 continuous functions descend through profinite quotient maps;
 finitely many $\Qlbar$-representatives admit one common finite field.
 \item \textbf{Record (do not define by) comparison identities:}
 the ordinary limit, rational localization, and finite-field colimit
 formulas from Proposition~\ref{prop:comparison-only}.
 \item \textbf{Build profinite test objects over schemes:}
 $X[S]$ using finite quotients of $S$, and construct
 $p_X:X_{\proet}\to*_{\proet}$.
 \item \textbf{Define pulled-back coefficients:}
 $A_X=p_X^{-1}\Lambda$; for the four local fields calculate
 $C(|U|,T)$, and for $\Qlbar$ use the locally finite-stage
 data of Proposition~\ref{prop:Qbar-X}.
 \item \textbf{Construct the coherent pullback maps:}
 $\alpha_f^\Lambda:f^{-1}A_Y\xrightarrow{\sim}A_X$,
 with unit, composition, and coefficient-change naturality.
 \item \textbf{Fix the lax orientation:}
 general structure maps are $\theta_f:f^{-1}A_Y\to A_X$;
 for coefficients pulled back from the point, all $\theta_f$
 are the invertible maps $\alpha_f^\Lambda$.
\end{enumerate}

\noindent\textbf{Conclusion.}
Every coefficient ring is specified concretely on the profinite point,
including $\Qlbar$ by \emph{finite-extension representatives} rather
than by a preliminary topology or a categorical colimit.  Passing to
$X_{\proet}$ is one uniform geometric inverse-image construction.
The resulting ring sheaves, together with
$f^{-1}A_Y\xrightarrow{\sim}A_X$, constitute a coherent Cartesian
coefficient system in the chosen pullback-oriented lax convention.

\begin{thebibliography}{9}
\bibitem{BS}
B.~Bhatt and P.~Scholze, \emph{The pro-\'{e}tale topology for schemes},
Ast\'{e}risque \textbf{369} (2015), 99--201, especially
Sections~4.1--4.3 and Lemmas~4.2.12, 4.3.7.
\url{https://arxiv.org/abs/1309.1198}.
\bibitem{HRS}
T.~Hemo, T.~Richarz and J.~Scholbach,
\emph{Constructible sheaves on schemes},
Advances in Mathematics \textbf{429} (2023), 109179,
especially Example~3.1 and Lemma~5.6.
\url{https://arxiv.org/abs/2305.18131}.
\bibitem{ScholzeCond}
P.~Scholze, \emph{Lectures on Condensed Mathematics},
particularly the definition of condensed sets and the
continuous-function construction.
\url{https://people.mpim-bonn.mpg.de/scholze/Condensed.pdf}.
\bibitem{GHN}
D.~Gepner, R.~Haugseng and T.~Nikolaus,
\emph{Lax colimits and free fibrations in $\infty$-categories},
Documenta Mathematica \textbf{22} (2017), 1225--1266.
For the distinction between sections of Cartesian fibrations and
(op)lax limits.
\end{thebibliography}
\end{document}
